\chapter{Axioms}
\label{chap:axioms}

\section{Minimal foundation}

The framework rests on four claims, the first three of which are essentially
established physics and the fourth of which is a minimal natural extension.

\begin{axiom}[Quantum mechanics applies]\label{ax:qm}
The supraverse is described by quantum field theory. Vacuum states have nonzero
fluctuations governed by the uncertainty principle. Over infinite spacetime,
fluctuations of arbitrary magnitude occur with probability one.
\end{axiom}

\noindent\emph{The timeless form of this axiom.} Axiom~\ref{ax:qm} is meant in
the timeless, quantum-gravitational sense. Textbook quantum mechanics presupposes
a background time---a Schr\"odinger equation with a $t$---but the supraverse as a
whole has no external clock (Chapter~\ref{chap:void}); time is emergent,
manufactured locally at each bounce. What we assume is therefore the
Wheeler--DeWitt form: the supraverse is a quantum \emph{state}, not a
process---a single stationary amplitude. Read this way, the ``fluctuations'' named
above are not events occurring in sequence but the timeless structure of that one
state---its amplitude having (exponentially small) support on configurations that
contain condensed universes (Chapter~\ref{chap:void},
Section~\ref{sec:void-eigenstate}). This is an extrapolation of tested physics
into a regime we cannot probe, and we label it as such; but it is the
\emph{conservative} extrapolation---it adds no new law, only the universality of
the law we already have.

\begin{axiom}[Special relativity applies]\label{ax:sr}
Spacetime has Lorentzian signature. Causal structure is determined by light
cones. CPT symmetry holds for any local Lorentz-invariant QFT.
\end{axiom}

\begin{axiom}[Einstein--Cartan gravity applies]\label{ax:ec}
The connection on spacetime is allowed to be asymmetric. The antisymmetric part
(torsion) couples to fermion spin via the axial current. In the absence of
fermionic matter, the theory reduces to standard general relativity. With
fermions, torsion produces an effective interaction between spins that scales as
the square of the spin density, and is therefore always positive (repulsive at
high spin density) regardless of relative spin orientation.
\end{axiom}

\noindent The Einstein--Cartan Lagrangian, in shorthand:
\[
  \mathcal{L}_{EC} = \frac{1}{2\kappa}(R - 2\Lambda)
  + \mathcal{L}_{\text{matter}}(\psi, \nabla\psi, g, T),
\]
where $R$ is the curvature scalar built from the full (torsion-including)
connection, and the covariant derivative $\nabla$ acts on fermionic fields
$\psi$ with torsion contributions.

\begin{axiom}[There exists a bounce density]\label{ax:rhoc}
At sufficiently high mass--energy density, torsion-mediated repulsion exceeds
gravitational attraction. The critical density (the bounce density) is
approximately
\[
  \rhoc \sim \frac{m_P^4}{\hbar^3}\,\alpha_{EC},
\]
where $m_P$ is the Planck mass and $\alpha_{EC}$ is a dimensionless
Einstein--Cartan coupling.
\end{axiom}

\noindent Estimates place this at
$\rhoc \sim 10^{50}\,\mathrm{kg\,m^{-3}}$, far below Planck density
($10^{96}\,\mathrm{kg\,m^{-3}}$) but well above nuclear density
($10^{17}\,\mathrm{kg\,m^{-3}}$). This is calculable from
Axioms~\ref{ax:qm}--\ref{ax:ec} in principle; Axiom~\ref{ax:rhoc} is the
explicit acknowledgment that this density exists and matters.

\section{What we do not assume}

\begin{itemize}
  \item No Big Bang as a singular event.
  \item No initial conditions for the universe.
  \item No fine-tuned cosmological constant or inflaton field.
  \item No special low-entropy past as a separate postulate.
  \item No additional dimensions beyond four.
  \item No additional forces or fields beyond Standard Model $+$ gravity $+$ torsion.
  \item No anthropic principle as a separate axiom (it emerges from the framework).
\end{itemize}

\section{What follows automatically}

From Axioms~\ref{ax:qm}--\ref{ax:rhoc} plus infinity (infinite spacetime
volume), the following are consequences rather than assumptions:

\begin{enumerate}
  \item \textbf{Spontaneous OG bounce nucleation.} Rare vacuum fluctuations reach
    bounce density. By Axiom~\ref{ax:qm} plus infinity, this happens somewhere
    with probability one. By Axiom~\ref{ax:rhoc}, such fluctuations bounce rather
    than collapsing to singularities.
  \item \textbf{Baby universe production.} The bounce geometry produces new
    spacetime regions (developed in Chapter~\ref{chap:dawn}).
  \item \textbf{CPT-symmetric supraverse.} Vacuum fluctuations are statistically
    symmetric in matter/antimatter content. Specific fluctuations have specific
    biases, but the supraverse-wide distribution is symmetric.
  \item \textbf{Thermodynamic equilibrium configuration.} Gravitational entropy
    considerations (high entropy $=$ clumped/black-hole-rich; low entropy $=$
    smooth/uniform) drive the supraverse toward a foam of universes as its
    equilibrium configuration.
  \item \textbf{Local arrow of time.} Each universe has a low-entropy bounce in
    its past and a high-entropy future, producing a local thermodynamic
    time-arrow.
  \item \textbf{Observer placement statistics.} Observers exist in viable
    universes near the peak of the supraverse population distribution.
  \item \textbf{Information preservation.} Hawking radiation entangles parent
    exterior with baby interior. The full supraverse is unitary even though
    individual universes appear to lose information.
\end{enumerate}

\section{On the role of infinity}

The framework requires the supraverse to be effectively infinite---at minimum,
large enough that rare fluctuations to bounce density occur. ``Effectively
infinite'' in this context means: large enough relative to the characteristic
timescales of vacuum fluctuation that the relevant fluctuations are
statistically guaranteed to occur.

For a finite supraverse, some fluctuations might never occur and the framework
would have to specify which ones do. This introduces an arbitrary cutoff that we
want to avoid. The natural assumption is true infinity, which removes this
arbitrariness at the cost of being philosophically heavier.

The supraverse being infinite does not require it to be temporally infinite in
some absolute sense---it requires only that enough ``supraverse time'' exist for
the relevant condensation to occur. Whether the supraverse has a beginning at
the supraverse level is a question we do not currently need to answer; we only
need it to be old enough that condensation has reached equilibrium.

\section{On the role of Einstein--Cartan}

Einstein--Cartan is the minimal natural extension of general relativity. It is
what you get by dropping the artificial restriction that the connection must be
symmetric. The connection's antisymmetric part (torsion) is mathematically
natural and couples to a physical quantity that exists in nature (fermion spin).
At ordinary densities, torsion effects are utterly negligible
($\sim 10^{-40}$ corrections to GR). At bounce density, they dominate.

We do not postulate Einstein--Cartan as a new theory; we recognize it as the
geometric framework that should have been used all along, with standard general
relativity as the low-density limit. This is not an additional assumption beyond
standard physics---it is the removal of an assumption (the symmetry of the
connection).

\section{The whole framework, in one breath}
\label{sec:epistemic-copernicus}

It is worth stating plainly what kind of claim this monograph makes, because the
reach of the conclusion can obscure how little it assumes. Every ingredient is
something we already measure here on Earth:
\begin{itemize}
  \item a \textbf{quantum vacuum that fluctuates}---measured, down to the Casimir
    force and the Lamb shift (Axiom~\ref{ax:qm});
  \item \textbf{gravity}---measured, and in its modern form geometric
    (Axiom~\ref{ax:sr}, Axiom~\ref{ax:ec});
  \item \textbf{torsion repulsion at extreme density}---not directly measured, but
    \emph{forced} on us the moment we let fermion spin source geometry, which it must
    (Axiom~\ref{ax:ec}, Axiom~\ref{ax:rhoc}).
\end{itemize}
From these three, plus infinity, a universe like our own is born in a never-stopping
loop: rare fluctuations reach bounce density, bounce instead of collapsing, extrude
baby universes, grow, heat-die into smooth voids, and seed the next generation---
forever, with no first cause and no fine-tuning (\Cref{chap:void}). We do not
have to \emph{believe} this happens; given what we already observe, it \emph{must}.

This is the epistemic inversion at the heart of the program---call it
\emph{epistemic Copernicus}. The Copernican principle says we should not assume we
occupy a special place; here that principle is not assumed but \emph{returned} as a
result, because the recursive supraverse makes our Big-Bang vantage a typical one
among uncountably many (\Cref{chap:dawn}). And the framework's central claim is
not ``the universe might be like this'' but the stronger, humbler ``\emph{from what
we have already observed, something like us must exist}.'' We are not arguing the
world into a shape; we are reading off what the measured world already entails. The
work of the rest of this book is to make that entailment precise, quantitative, and
falsifiable---and, where a more likely mechanism might exist that we have not yet
found, to say so honestly. If a better one is found, it wins; until then, this is
enough to require our own existence.
